\documentclass[11pt,a4paper]{article}
\usepackage[utf8]{inputenc}
\usepackage[T1]{fontenc}
\usepackage[spanish,es-tabla]{babel}
\usepackage{lmodern}
\usepackage[top=2.5cm,bottom=2.5cm,left=2.5cm,right=2.5cm,headheight=14pt]{geometry}
\usepackage{microtype}
\usepackage[table,xcdraw]{xcolor}
\definecolor{headerblue}{RGB}{24,43,73}
\definecolor{accentblue}{RGB}{30,80,150}
\definecolor{lightbg}{RGB}{248,249,251}
\definecolor{bordergray}{RGB}{220,224,230}
\definecolor{darktext}{RGB}{33,37,41}
\definecolor{codebg}{RGB}{245,247,250}
\definecolor{commentgreen}{RGB}{40,140,70}
\definecolor{keywordpurple}{RGB}{120,40,140}
\definecolor{stringblue}{RGB}{20,90,160}
\colorlet{bannerblue}{headerblue}
\usepackage{amsmath,amssymb}
\usepackage{graphicx}
\graphicspath{{assets/hijos/}{assets/linus/}{assets/sinergia/}{assets/modus/}{assets/}{../assets/}{../assets/hijos/}{../assets/linus/}{../assets/sinergia/}{../assets/modus/}}
\usepackage{hyperref}
\hypersetup{colorlinks=true,linkcolor=accentblue,citecolor=accentblue,urlcolor=accentblue,bookmarksnumbered=true,breaklinks=true}
\usepackage{booktabs,tabularx,longtable,array,enumitem,caption,float}
\usepackage{fancyhdr}
\pagestyle{fancy}
\fancyhf{}
\renewcommand{\headrulewidth}{0.4pt}
\fancyhead[L]{\small\textsc{LogiLearn}}
\fancyhead[R]{\small\leftmark}
\fancyfoot[C]{\small\thepage}
\usepackage{listings}
\lstdefinelanguage{TypeScript}{keywords={const,let,function,return,import,export,for,if,else,new},keywordstyle=\color{accentblue}\bfseries,comment=[l]{//},commentstyle=\color{gray}\itshape,sensitive=true}
\lstset{basicstyle=\ttfamily\footnotesize,breaklines=true,frame=single,numbers=left,numberstyle=\tiny\color{gray},showstringspaces=false,tabsize=2,extendedchars=true,literate={á}{{\'a}}1 {é}{{\'e}}1 {í}{{\'i}}1 {ó}{{\'o}}1 {ú}{{\'u}}1 {ñ}{{\~n}}1 {Á}{{\'A}}1 {É}{{\'E}}1 {Í}{{\'I}}1 {Ó}{{\'O}}1 {Ú}{{\'U}}1 {Ñ}{{\~N}}1 {∧}{{$\land$}}1 {∨}{{$\lor$}}1 {¬}{{$\neg$}}1 {→}{{$\to$}}1 {↔}{{$\leftrightarrow$}}1 {∀}{{$\forall$}}1 {∃}{{$\exists$}}1 {≡}{{$\equiv$}}1 {✓}{{$\checkmark$}}1}
\usepackage{tcolorbox}
\begin{document}
\begin{titlepage}
\centering
{\Large Universidad Nacional Pedro Ruiz Gallo\par}
{\large Escuela Profesional de Ingenier\'ia de Sistemas --- MATG1001\par}
\vspace{0.6cm}
\includegraphics[width=\textwidth]{hero_cover.png}\par
\vspace{0.6cm}
{\Huge \bfseries\color{headerblue} Manual de Uso\par}
\vspace{0.2cm}
{\Large LogiLearn --- Gu\'ia paso a paso por equipo\par}
\vspace{0.4cm}
{\large Plataforma LogiLearn --- L\'ogica Simb\'olica Global\par}
{\small Docente: Dr. Mardo Victor Gonzales Herrera --- Septiembre 2026\par}
\end{titlepage}
\tableofcontents
\newpage

\section{Cap\'itulo 1 --- Sinergia: Tablas de Verdad}

\subsection{Acceso al Módulo}
- En la nube: Ingrese a la URL de GitHub Pages y haga clic en la pestaña \textbf{«Tablas de verdad»}.
- En local: Ejecute \texttt{npm run dev} y navegue a \texttt{http://localhost:5173/tablas}.

\subsection{Notaciones Aceptadas}
\begin{table}[H]
\centering
\caption{Notaciones aceptadas por el motor lógico de LogiLearn}
\label{tab:notaciones}
\small
\begin{tabular}{lccc}
\toprule
\textbf{Conectivo} & \textbf{Símbolo Unicode} & \textbf{ASCII} & \textbf{Palabras clave} \\
\midrule
Negación \& $\neg$ \& \texttt{\~{}} , \texttt{!} , \texttt{-} \& \texttt{NOT} \\
Conjunción \& $\land$ \& \texttt{\^{}} , \texttt{\&} , \texttt{.} \& \texttt{AND} \\
Disyunción \& $\lor$ \& \texttt{|} , \texttt{+} \& \texttt{OR} \\
Implicación \& $\to$ \& \texttt{->} , \texttt{=>} \& \texttt{IMPLIES} \\
Bicondicional \& $\leftrightarrow$ \& \texttt{<->} , \texttt{<=>} \& \texttt{IFF} \\
\bottomrule
\end{tabular}
\end{table}

\subsection{Interpretación de la Tabla y Clasificación}
\begin{itemize}[leftmargin=*]
  \item \textbf{Columnas de Variables:} Asignaciones $V$ y $F$ para cada proposición simple ($P, Q, R$).
  \item \textbf{Columnas de Subexpresiones:} Evaluación intermedia de los paréntesis de adentro hacia afuera.
  \item \textbf{Columna Resultado:} Matriz principal clasificada como \textbf{TAUTOLOGÍA} (todo V), \textbf{CONTRADICCIÓN} (todo F) o \textbf{CONTINGENCIA} (mixta).
\end{itemize}

\% ==============================================================================
\% 6. REGISTRO DE EVOLUCIÓN POR SPRINTS
\% ==============================================================================
\section{Registro de Evolución por Sprints}
8. DOCUMENTACIÓN DE LAS PESTAÑAS COMPLEMENTARIAS
\begin{lstlisting}[language=TypeScript,basicstyle=\ttfamily\scriptsize]
================================================================================
\end{lstlisting}
\subsection{8.1 Pestaña Inicio (HomeView): Punto de entrada, presentación y enlaces de navegación (src/pages/Home.vue).}
\subsection{8.2 Pestaña Aprender (LearnView): Recorrido guiado en stepper de 4 pasos (Definición, Ejemplo, Ejercicio, Solución), banner de repaso inteligente y botón ''Practicar ahora'' (src/pages/aprender/index.vue).}
\subsection{8.3 Pestaña Leyes Lógicas (LogicLawsView): Biblioteca interactiva de las 12 leyes lógicas con interfaz LeyLogica y buscador reactivo en tiempo real (src/pages/leyes-logicas/index.vue).}
\subsection{8.4 Pestaña Ejercicios (ExercisesView / ExerciseRunnerView): Catálogo completo de práctica interactiva con doble filtro combinable por dificultad y tema, sincronización con query param ?tema= y feedback con OptionPill (src/pages/ejercicios/).}
\subsection{8.5 Pestaña Progreso (ProgressView): Historial y estadísticas del usuario persistidas en localStorage con reactive() de Vue 3, sistema de logros reactivo y panel de progreso por tema (src/pages/progreso/index.vue).}
\subsection{8.6 Flujo de Integración entre Pestañas (Figura A): Ciclo de aprendizaje cerrado articulado por el store de progreso.}

\begin{lstlisting}[language=TypeScript,basicstyle=\ttfamily\scriptsize]
================================================================================
\end{lstlisting}
\begin{figure}[H]\centering
\includegraphics[width=0.9\textwidth]{01_encabezado_navbar.jpg}
\caption{Evidencia del m'odulo (Equipo Sinergia)}
\end{figure}
\begin{figure}[H]\centering
\includegraphics[width=0.9\textwidth]{02_modulo_tablas_verdad.jpg}
\caption{Evidencia del m'odulo (Equipo Sinergia)}
\end{figure}
\begin{figure}[H]\centering
\includegraphics[width=0.9\textwidth]{extracted_img_02.jpg}
\caption{Evidencia del m'odulo (Equipo Sinergia)}
\end{figure}
\begin{figure}[H]\centering
\includegraphics[width=0.9\textwidth]{extracted_img_03.jpg}
\caption{Evidencia del m'odulo (Equipo Sinergia)}
\end{figure}
\begin{figure}[H]\centering
\includegraphics[width=0.9\textwidth]{extracted_img_04.jpg}
\caption{Evidencia del m'odulo (Equipo Sinergia)}
\end{figure}
\begin{figure}[H]\centering
\includegraphics[width=0.9\textwidth]{extracted_img_05.jpg}
\caption{Evidencia del m'odulo (Equipo Sinergia)}
\end{figure}
\begin{figure}[H]\centering
\includegraphics[width=0.9\textwidth]{extracted_img_06.jpg}
\caption{Evidencia del m'odulo (Equipo Sinergia)}
\end{figure}
\section{Cap\'itulo 2 --- Hijos de Linus: Inferencias}
\section{Manual de Uso --- Módulo de Validez de Inferencias Lógicas}

\textbf{Equipo:} Hijos de Linus
\textbf{Versión:} 1.3
\textbf{Fecha:} Septiembre 2026

\hrulefill

\subsubsection{Tabla de Contenidos}

1. ¿Qué es este módulo? (\url{#1-qué-es-este-módulo})
2. ¿Para qué sirve? (\url{#2-para-qué-sirve})
3. ¿Cómo se accede? (\url{#3-cómo-se-accede})
4. Pantalla principal: descripción general (\url{#4-pantalla-principal-descripción-general})
5. Pestaña 1 --- Simbología Formal (\url{#5-pestaña-1--simbología-formal})
6. Teclado Simbólico (\url{#6-teclado-simbólico})
7. Ejemplos Rápidos (\url{#7-ejemplos-rápidos})
8. Pestaña 2 --- Lenguaje Natural (\url{#8-pestaña-2--lenguaje-natural})
9. Pestaña 3 --- Árbol Sintáctico (AST) (\url{#9-pestaña-3--árbol-sintáctico-ast})
10. Como interpretar el diagrama de Árbol (\url{#10-como-interpretar-el-diagrama-de-árbol})
11. Resultado de la Inferencia (\url{#11-resultado-de-la-inferencia})
12. Panel de Trazabilidad (Paso a Paso) (\url{#12-panel-de-trazabilidad-paso-a-paso})
13. Exportación Académica (Markdown y LaTeX) (\url{#13-exportación-académica-markdown-y-latex})
14. Historial de Ejercicios (\url{#14-historial-de-ejercicios})
15. Los 8 Operadores Lógicos: Referencia Completa (\url{#15-los-8-operadores-lógicos-referencia-completa})
16. Las 11 Reglas de Inferencia (\url{#16-las-11-reglas-de-inferencia})
17. Preguntas Frecuentes (FAQ) (\url{#17-preguntas-frecuentes-faq})
18. Glosario de Términos (\url{#18-glosario-de-términos})

\hrulefill

\subsection{1. ¿Qué es este módulo?}

El módulo \textbf{'Validez de Inferencias Lógicas'} es una herramienta interactiva desarrollada por el equipo \textbf{Hijos de Linus} que permite:

\begin{itemize}[leftmargin=*]
\item Escribir argumentos lógicos (premisas + conclusión) en un formulario.
\item Verificar si la conclusión se deduce válidamente de las premisas.
\item Visualizar la estructura interna del argumento como un árbol jerárquico.
\item Traducir los símbolos formales a lenguaje natural (español).
\item Obtener una demostración paso a paso con las reglas lógicas aplicadas.
\item Exportar la demostración en formato académico (Markdown o LaTeX).
\end{itemize}

En palabras simples: \textbf{es una calculadora de lógica proposicional} que te dice si tu argumento es correcto o no, y te explica por qué paso a paso.

\hrulefill

\subsection{2. ¿Para qué sirve?}

Uso --- Ejemplo
\textbf{Estudiar lógica proposicional} --- Verificar si un argumento es válido antes de entregar una tarea
\textbf{Aprender reglas de inferencia} --- Entender qué regla se aplica en cada paso de una demostración
\textbf{Visualizar la estructura de fórmulas} --- Ver cómo se descompone jerárquicamente una expresión lógica
\textbf{Traducir símbolos a texto} --- Convertir fórmulas como \texttt{P  Q  R} a oraciones en español
\textbf{Generar demostraciones formales} --- Exportar un paso a paso en LaTeX para un trabajo académico
\textbf{Detectar falacias lógicas} --- Saber por qué un argumento inválido falla y cómo corregirlo

\hrulefill

\subsection{3. ¿Cómo se accede?}

\subsubsection{Paso 1: Iniciar el servidor de desarrollo}

Si eres desarrollador, abre una terminal en la carpeta del proyecto y ejecuta:

\begin{lstlisting}[language={},basicstyle=\ttfamily\scriptsize]
npm install
npm run dev
\end{lstlisting}

El servidor arrancará en \texttt{http://localhost:5173}.

\subsubsection{Paso 2: Abrir el navegador}

Abre tu navegador web (Chrome, Firefox, Edge, Safari) y ve a:

\begin{lstlisting}[language={},basicstyle=\ttfamily\scriptsize]
http://localhost:5173
\end{lstlisting}

\subsubsection{Paso 3: Navegar al módulo}

Desde la página principal, haz clic en la sección o tarjeta de \textbf{'Inferencias Lógicas'} (o navega directamente a la ruta \texttt{/inferencias}).

\hrulefill

\subsection{4. Pantalla principal: descripción general}

Al abrir el módulo verás la siguiente estructura:

\begin{lstlisting}[language={},basicstyle=\ttfamily\scriptsize]
  [Volver] [Hijos de Linus]              [Historial (N)] 
                                                         
   Validez de Inferencias Logicas                       
    Ingresa premisas formales, visualiza el arbol...     

                                                         
  [Simbologia Formal] [Lenguaje Natural] [Arbol AST]     
                                                         
       
                           Resultado de la        |   
    Contenido de la        Inferencia             |   
    pestana activa                                |   
                           Trazabilidad paso a    |   
                           paso con explicaciones |   
                           y exportacion          |   
       |
                                                         
\end{lstlisting}

\textbf{Zona izquierda:} Contiene las 3 pestañas para interactuar con el argumento.
\textbf{Zona derecha:} Muestra el resultado (válida/inválida) y la demostración paso a paso.

\hrulefill

\subsection{5. Pestaña 1 --- Simbología Formal}

Esta es la pestaña principal donde \textbf{escribes tu argumento lógico}.

\subsubsection{5.1 Escribir Premisas}

Las \textbf{premisas} son las afirmaciones de las partes de las cuales quieres deducir algo. Cada premisa va en una línea separada.

\textbf{Ejemplo:} Si quieres expresar 'Si llueve, la calle se moja' y 'Llueve', escribes:

\begin{lstlisting}[language={},basicstyle=\ttfamily\scriptsize]
P ENTONCES Q
P
\end{lstlisting}

Esto significa:
\begin{itemize}[leftmargin=*]
\item Línea 1: Si P, entonces Q (P $ \to $ Q)
\item Línea 2: P es verdadero
\end{itemize}

> \textbf{Nota:} El campo de premisas muestra números de línea a la izquierda (1, 2, 3...) para facilitar la referencia.

\subsubsection{5.2 Escribir la Conclusión}

La \textbf{conclusión} es lo que quieres demostrar que se deduce de las premisas. Va en un solo campo debajo de las premisas.

\textbf{Ejemplo:} Para concluir que 'Q' (la calle se moja):

\begin{lstlisting}[language={},basicstyle=\ttfamily\scriptsize]
 Q
\end{lstlisting}

El símbolo \texttt{} (por lo tanto) aparece automáticamente a la izquierda del campo.

\subsubsection{5.3 Dos formas de escribir}

Puedes escribir de \textbf{dos maneras equivalentes}:

Forma matemática --- Forma con palabras clave
\texttt{P  Q} --- \texttt{P ENTONCES Q}
\texttt{P  Q} --- \texttt{P Y Q}
\texttt{P  Q} --- \texttt{P O Q}
\texttt{P} --- \texttt{NO P}
\texttt{P  Q} --- \texttt{P SI\_Y\_SOLO\_SI Q}

El sistema acepta \textbf{ambas formas}. Si escribes símbolos matemáticos ($ \to $, $ \land $, $ \lor $, $ \lnot $, $ \leftrightarrow $, , $ \downarrow $, $ \uparrow $), el sistema los convierte automáticamente al formato interno.

\subsubsection{5.4 Botón 'Demostrar Inferencia'}

Una vez escrito el argumento, haz clic en el botón negro \textbf{'Demostrar Inferencia $ \to $'} en la parte inferior del formulario. El sistema procesará tu argumento y mostrará el resultado en la columna derecha.

\hrulefill

\subsection{6. Teclado Simbólico}

Debajo del formulario de entrada encontrarás un \textbf{teclado virtual} con botones para insertar símbolos. Esto es especialmente útil si no tienes símbolos en tu teclado.

\subsubsection{6.1 Fila superior --- Conectivos lógicos}

Botón --- Símbolo --- Nombre --- Ejemplo de uso
\texttt{} --- $ \lnot $ --- Negación (NO) --- \texttt{P} = 'No P'
\texttt{} --- $ \land $ --- Conjunción (Y) --- \texttt{P  Q} = 'P y Q'
\texttt{} --- $ \lor $ --- Disyunción (O) --- \texttt{P  Q} = 'P o Q'
\texttt{} --- Disyunción Exclusiva (XOR) --- \texttt{P  Q} = 'P o Q, pero no ambos'
\texttt{} --- $ \to $ --- Implicación (ENTONCES) --- \texttt{P  Q} = 'Si P entonces Q'
\texttt{} --- $ \leftrightarrow $ --- Bicondicional (SI Y SOLO SI) --- \texttt{P  Q} = 'P si y solo si Q'
\texttt{(} --- ( --- Paréntesis de apertura --- Para agrupar
\texttt{)} --- ) --- Paréntesis de cierre --- Para agrupar

\subsubsection{6.2 Fila inferior --- Variables y acciones}

Botón --- Función
\texttt{P}, \texttt{Q}, \texttt{R}, \texttt{S} --- Variables proposicionales (grupo 1)
\texttt{A}, \texttt{B}, \texttt{C}, \texttt{D} --- Variables proposicionales (grupo 2)
\texttt{Borrar} --- Borra el carácter anterior al cursor
\texttt{Salto} --- Inserta un salto de línea (solo en premisas)

\subsubsection{6.3 Cómo funciona la inserción}

\begin{itemize}[leftmargin=*]
\item Haz clic en el campo de \textbf{premisas} o \textbf{conclusión} donde quieras escribir.
\item El cursor se posicionará en ese campo.
\item Haz clic en cualquier botón del teclado virtual: el símbolo se insertará en la posición del cursor.
\item Los conectivos binarios ($ \land $, $ \lor $, , $ \to $, $ \leftrightarrow $) se insertan \textbf{con espacios automáticos} alrededor para mantener la legibilidad.
\end{itemize}

\hrulefill

\subsection{7. Ejemplos Rápidos}

En la parte superior del formulario encontrarás \textbf{4 ejemplos predefinidos} que puedes cargar con un solo clic:

\subsubsection{Ejemplos válidos (de color verde)}

Nombre --- Premisas --- Conclusión --- Tipo
\textbf{Modus Ponens} --- P $ \to $ Q, P --- Q --- Válido
\textbf{Silogismo Hipotético} --- P $ \to $ Q, Q $ \to $ R --- P $ \to $ R --- Válido
\textbf{Dilema Constructivo} --- P $ \to $ Q, R $ \to $ S, P $ \lor $ R --- Q $ \lor $ S --- Válido

\subsubsection{Ejemplo inválido (de color ámbar)}

Nombre --- Premisas --- Conclusión --- Tipo
\textbf{Falacia Afirm. Consecuente} --- P $ \to $ Q, Q --- P --- Falacia

\textbf{Para usar un ejemplo:} Haz clic en el botón del ejemplo. Los campos de premisas y conclusión se llenarán automáticamente. Luego haz clic en 'Demostrar Inferencia'.

\hrulefill

\subsection{8. Pestaña 2 --- Lenguaje Natural}

Esta pestaña \textbf{traduce tus fórmulas lógicas a oraciones en español} para facilitar la comprensión.

\subsubsection{8.1 Asignar significado a las variables}

En la parte superior verás tarjetas con cada variable detectada (P, Q, R, S, etc.). Cada una tiene un campo editable donde puedes escribir qué significa esa variable en contexto real.

\textbf{Valores por defecto:}

Variable --- Significado por defecto
P --- llueve
Q --- la calle se moja
R --- el suelo está resbaloso
S --- hay tráfico

\textbf{Para personalizar:} Haz clic en el campo de texto junto a la variable y escribe el significado que desees. Por ejemplo, podrías cambiar P a 'estudia mucho' y Q a 'aprueba el examen'.

\subsubsection{8.2 Ver la traducción}

Una vez asignados los significados, en la parte inferior aparecerá el \textbf{'Argumento Traducido en Prosa'}:

\textbf{Ejemplo} con las premisas \texttt{P  Q} y \texttt{P}, y conclusión \texttt{Q}:

> \textbf{Premisa 1:} Si llueve, entonces la calle se moja.
> \textbf{Premisa 2:} Llueve.
> \textbf{$ \therefore $ Conclusión:} Por lo tanto, la calle se moja.

Esto te permite leer el argumento como una oración normal en español, sin símbolos.

\hrulefill

\subsection{9. Pestaña 3 --- Árbol Sintáctico (AST)}

Esta pestaña muestra la \textbf{estructura jerárquica} de tu argumento como un árbol visual. Es la función principal implementada por \textbf{Hijos de Linus}.

\subsubsection{9.1 ¿Qué es un AST?}

\textbf{AST} significa \textbf{Árbol de Sintaxis Abstracta} (Abstract Syntax Tree). Es una representación visual que muestra cómo se 'descompone' una fórmula lógica en sus partes componentes.

Piensa en ello como un \textbf{árbol genealógico de la fórmula}: la raíz es la operación principal, y las ramas son las sub-fórmulas que componen cada parte.

\subsubsection{9.2 Descripción de los elementos del árbol}

\paragraph{Nodo raíz: 'Demostración'}
\begin{itemize}[leftmargin=*]
\item Es el nodo superior oscuro con icono de red.
\item Representa la deducción completa (todas las premisas + conclusión).
\item De él salen ramas hacia cada premisa y la conclusión.
\end{itemize}

\paragraph{Ramas principales}
\begin{itemize}[leftmargin=*]
\item Cada rama principal representa una \textbf{proposición individual}:
\item \texttt{P1}, \texttt{P2}, etc. = Premisas
\item \texttt{ Conclusion} = La conclusión
\end{itemize}

\paragraph{Nodos de operador}
\begin{itemize}[leftmargin=*]
\item Son los rectángulos coloreados con un símbolo ($ \lnot $, $ \land $, $ \lor $, $ \to $, etc.).
\item Cada color corresponde a un tipo de operador.
\item Si un nodo tiene \textbf{dos hijos}, es un operador binario (como $ \land $ o $ \to $).
\item Si tiene \textbf{un solo hijo}, es un operador unario (como $ \lnot $).
\end{itemize}

\paragraph{Hojas (variables)}
\begin{itemize}[leftmargin=*]
\item Son los nodos grises al final de las ramas.
\item Representan las variables proposicionales (P, Q, R, etc.).
\item No tienen hijos: son los 'elementos más pequeños' de la fórmula.
\end{itemize}

\subsubsection{9.3 Colores de los operadores}

Color --- Operador --- Símbolo --- Ejemplo
Rojo --- Negación --- $ \lnot $ --- \texttt{P}
Azul --- Conjunción --- $ \land $ --- \texttt{P  Q}
Esmeralda --- Disyunción --- $ \lor $ --- \texttt{P  Q}
Ámbar --- Disyunción Exclusiva --- \texttt{P  Q}
Índigo --- Implicación --- $ \to $ --- \texttt{P  Q}
Cian --- Bicondicional --- $ \leftrightarrow $ --- \texttt{P  Q}
Fucsia --- Nor --- $ \downarrow $ --- \texttt{P  Q}
Naranja --- Nand --- $ \uparrow $ --- \texttt{P  Q}

\subsubsection{9.4 Leyenda de operadores}

En la parte superior del árbol verás una \textbf{leyenda} con todos los operadores y sus colores. Esto te ayuda a identificar rápidamente qué tipo de conexión representa cada nodo.

\subsubsection{9.5 Guía de interpretación}

Debajo de la leyenda hay un panel azul claro con la \textbf{'Interpretación del árbol'} que explica:

\begin{itemize}[leftmargin=*]
\item \textbf{Nodo raíz:} representa la deducción completa.
\item \textbf{Ramas principales:} premisas individuales (P1, P2...) y conclusión ($ \therefore $).
\item \textbf{Nodos y hojas:} operadores lógicos como conectores y letras como variables.
\end{itemize}

\subsubsection{9.6 Estadísticas del árbol}

Debajo de la leyenda verás \textbf{etiquetas con estadísticas} del árbol actual:

Estadística --- Significado
\textbf{N proposiciones} --- Número total de premisas + conclusión
\textbf{N nodos} --- Cantidad total de elementos en el árbol
\textbf{prof. máx N} --- Profundidad máxima (cuántos niveles tiene el árbol)
\textbf{N variables} --- Cantidad de letras variables (P, Q, R, etc.)
\textbf{N conectivos} --- Cantidad de operadores lógicos

\hrulefill

\subsection{10. Como interpretar el diagrama de Árbol}

\subsubsection{Ejemplo: \texttt{P  Q}, \texttt{P}  \texttt{Q}}

El árbol se ve así (esquema simplificado):

\begin{lstlisting}[language={},basicstyle=\ttfamily\scriptsize]
              Demostracion
            /      |       \
          P1      P2      Conclusion
          |       |         |
                 P         Q
         / \
        P   Q
\end{lstlisting}

\textbf{Explicación:}

1. \textbf{Raíz 'Demostración'} (nodo oscuro arriba): representa todo el argumento.
2. \textbf{Rama P1} (Premisa 1: \texttt{P  Q}):
\begin{itemize}[leftmargin=*]
\item Nodo \texttt{} (índigo): el operador principal es 'ENTONCES'.
\item Hijos: \texttt{P} (izquierda) y \texttt{Q} (derecha).
\end{itemize}
3. \textbf{Rama P2} (Premisa 2: \texttt{P}):
\begin{itemize}[leftmargin=*]
\item Nodo \texttt{P} (gris): es una variable directa, sin operador.
\end{itemize}
4. \textbf{Rama $ \therefore $ Conclusión} (Conclusión: \texttt{Q}):
\begin{itemize}[leftmargin=*]
\item Nodo \texttt{Q} (gris): es una variable directa.
\end{itemize}

\subsubsection{Reglas de lectura}

\begin{itemize}[leftmargin=*]
\item \textbf{De arriba a abajo:} desde la operación principal hasta las variables base.
\item \textbf{De izquierda a derecha:} en operadores binarios, el hijo izquierdo es la primera parte y el derecho la segunda.
\item \textbf{Colores:} ayudan a identificar el tipo de operación rápidamente.
\item \textbf{Líneas:} conectan cada nodo padre con sus hijos directos.
\end{itemize}

\hrulefill

\subsection{11. Resultado de la Inferencia}

Después de hacer clic en 'Demostrar Inferencia', en la columna derecha aparecerá un indicador con el resultado. Hay \textbf{4 posibles estados}:

\subsubsection{11.1 Inferencia válida (Demostrada)}

\begin{itemize}[leftmargin=*]
\item \textbf{Color:} Verde
\item \textbf{Icono:}  (check)
\item \textbf{Significado:} La conclusión se deduce correctamente de las premisas. El argumento es formalmente válido.
\item \textbf{Acción:} Puedes ver la demostración paso a paso en el panel de Trazabilidad.
\end{itemize}

\subsubsection{11.2 Inferencia inválida (Refutada)}

\begin{itemize}[leftmargin=*]
\item \textbf{Color:} Rojo
\item \textbf{Icono:}  (cruz)
\item \textbf{Significado:} El argumento \textbf{no es válido}. Se encontró un contraejemplo que refuta la inferencia.
\item \textbf{Acción:} El panel de Trazabilidad mostrará el \textbf{diagnóstico del fallo} con:
\item Una descripción del error lógico.
\item Una explicación de por qué falla el razonamiento.
\item Un \textbf{contraejemplo} con valores de verdad (V/F) para cada variable.
\item Una \textbf{sugerencia de corrección}.
\end{itemize}

\subsubsection{11.3 Inferencia válida (Método indirecto requerido)}

\begin{itemize}[leftmargin=*]
\item \textbf{Color:} Índigo (azul oscuro)
\item \textbf{Icono:} Brújula
\item \textbf{Significado:} El argumento es válido, pero no se pudo demostrar por métodos directos (forward chaining). Requiere métodos como \textbf{Reducción al absurdo} o \textbf{Prueba condicional}.
\item \textbf{Acción:} No hay demostración directa disponible, pero el diagnóstico indica que el argumento es válido.
\end{itemize}

\subsubsection{11.4 Error de sintaxis o procesamiento}

\begin{itemize}[leftmargin=*]
\item \textbf{Color:} Ámbar (amarillo)
\item \textbf{Icono:}  (triángulo de advertencia)
\item \textbf{Significado:} Hubo un problema al evaluar la expresión. Puede ser una fórmula mal escrita, paréntesis desbalanceados, o un operador no reconocido.
\item \textbf{Acción:} Revisa la sintaxis de tus fórmulas. Asegúrate de que los paréntesis estén balanceados y que uses operadores válidos.
\end{itemize}

\hrulefill

\subsection{12. Panel de Trazabilidad (Paso a Paso)}

Debajo del indicador de resultado se encuentra el \textbf{Panel de Trazabilidad}, que es el corazón de la demostración.

\subsubsection{12.1 Cuando la inferencia es válida}

El panel muestra una \textbf{lista de pasos numerados} que forman la demostración formal. Cada paso contiene:

\paragraph{Estructura de cada paso:}

\begin{lstlisting}[language={},basicstyle=\ttfamily\scriptsize]
  [Numero de paso]  Regla Aplicada           
                    [Linea X] [Linea Y]      
                                             
   Formula deducida                         
                                             
  [Como se deduce?]                        
\end{lstlisting}

\begin{itemize}[leftmargin=*]
\item \textbf{Número de paso:} La línea global en la demostración (premisas + pasos comparten numeración).
\item \textbf{Regla aplicada:} El nombre de la regla lógica usada (ej. 'Modus Ponendo Ponens').
\item \textbf{Líneas referenciadas:} Las líneas base de las cuales se deduce este paso.
\item \textbf{Fórmula deducida:} La expresión resultante de aplicar la regla.
\item \textbf{Botón '¿Cómo se deduce?':} Al hacer clic, se despliega una explicación detallada.
\end{itemize}

\paragraph{Contenido del acordeón desplegado:}

Al hacer clic en \textbf{'¿Cómo se deduce?'} de un paso, se muestran 3 secciones:

1. \textbf{Premisas base utilizadas:} Las líneas originales que se usaron para deducir este paso, con su expresión y rol.
2. \textbf{Regla aplicada:} El nombre completo de la regla, su alias, y una explicación de cómo funciona.
3. \textbf{Resultado:} La conclusión obtenida de aplicar la regla.

\subsubsection{12.2 Cuando la inferencia es inválida}

En lugar de pasos de demostración, el panel muestra un \textbf{diagnóstico formal} con 4 secciones:

\paragraph{1. Análisis del Problema}
Descripción del error detectado en el argumento.

\paragraph{2. ¿Por qué falla este razonamiento?}
Explicación detallada de la falacia o error lógico.

\paragraph{3. Contraejemplo que refuta la validez}
Un conjunto de valores de verdad (V o F) para cada variable que:
\begin{itemize}[leftmargin=*]
\item Hace \textbf{todas las premisas verdaderas}.
\item Pero la \textbf{conclusión resulta falsa}.
\end{itemize}

Esto prueba formalmente que el argumento no es válido.

\textbf{Ejemplo de contraejemplo:}

Variable --- Valor
P --- Falso (F)
Q --- Verdadero (V)

Con esta asignación, todas las premisas son V pero la conclusión es F $ \to $ el argumento es inválido.

\paragraph{4. Sugerencia de corrección}
Una propuesta de cómo modificar el argumento para que sea válido.

\hrulefill

\subsection{13. Exportación Académica (Markdown y LaTeX)}

Cuando la inferencia es válida, en la parte superior del panel de Trazabilidad aparece una \textbf{barra de exportación} con 4 botones:

\subsubsection{13.1 Copiar Markdown (Copiar MD)}

Copia al portapapeles la demostración en formato Markdown con símbolos Unicode legibles. Ideal para:
\begin{itemize}[leftmargin=*]
\item Pegar en un documento de Google Docs o Word.
\item Compartir en foros o chats.
\item Crear notas en Obsidian, Notion, etc.
\end{itemize}

\subsubsection{13.2 Descargar Markdown (.md)}

Descarga un archivo \texttt{demostracion\_logica.md} con la demostración completa en formato Markdown.

\subsubsection{13.3 Copiar LaTeX (Copiar LaTeX)}

Copia al portapapeles un \textbf{documento LaTeX completo y autocompilable}. Incluye:
\begin{itemize}[leftmargin=*]
\item Preámbulo estándar (\texttt{\textbackslash{}documentclass}, \texttt{amsmath}, \texttt{amssymb}, \texttt{babel} en español).
\item Sección de argumento con premisas numeradas.
\item Demostración en entorno \texttt{align*} con numeración global.
\item Símbolos convertidos a comandos LaTeX (ej. \texttt{\textbackslash{}rightarrow}, \texttt{\textbackslash{}land}, \texttt{\textbackslash{}neg}).
\end{itemize}

\subsubsection{13.4 Descargar LaTeX (.tex)}

Descarga un archivo \texttt{demostracion\_logica.tex} listo para compilar en:
\begin{itemize}[leftmargin=*]
\item \textbf{Overleaf} (editor online de LaTeX, la forma más fácil).
\item \textbf{TeXstudio}, \textbf{TeXShop}, o cualquier compilador local.
\item Soporta pdfLaTeX, XeLaTeX y LuaLaTeX.
\end{itemize}

\subsubsection{Formato de la exportación Markdown}

\begin{lstlisting}[language={},basicstyle=\ttfamily\scriptsize]
### Demostracion Formal de Inferencia Logica

**Premisas:**

1. P  Q
2. P

**Conclusion:**  Q

**Deduccion formal paso a paso:**

3. Q *[Modus Ponendo Ponens (Linea 1, Linea 2)]*
\end{lstlisting}

\subsubsection{Formato de la exportación LaTeX}

\begin{lstlisting}[language={},basicstyle=\ttfamily\scriptsize]
\documentclass[11pt]{article}
\usepackage[utf8]{inputenc}
\usepackage[T1]{fontenc}
\usepackage[spanish,es-noshorthands]{babel}
\usepackage{amsmath}
\usepackage{amssymb}
\usepackage[a4paper,margin=2.5cm]{geometry}

\title{Demostracion Formal de Inferencia Logica}

\begin{document}
\maketitle

\section*{Argumento}
\begin{enumerate}
  \item $P \rightarrow Q$
  \item $P$
\end{enumerate}

\section*{Demostracion formal}
\begin{align*}
(1) \quad & P \rightarrow Q && \text{Premisa} \\
(2) \quad & P && \text{Premisa} \\
\therefore (3) \quad & Q && \text{[Modus Ponendo Ponens (1, 2)]}
\end{align*}

\end{document}
\end{lstlisting}

\hrulefill

\subsection{14. Historial de Ejercicios}

En la esquina superior derecha del encabezado hay un botón \textbf{'Historial (N)'} que muestra los últimos 6 ejercicios que has resuelto.

\subsubsection{14.1 Cómo funciona}

\begin{itemize}[leftmargin=*]
\item Cada vez que haces clic en 'Demostrar Inferencia', el ejercicio se guarda automáticamente.
\item Se almacenan solo los \textbf{últimos 6} ejercicios (los más recientes reemplazan a los antiguos).
\item Los ejercicios duplicados se eliminan automáticamente.
\end{itemize}

\subsubsection{14.2 Cómo restaurar un ejercicio}

Haz clic en la tarjeta del ejercicio en el historial. Los campos de premisas y conclusión se llenarán automáticamente con los valores anteriores, y se procesará la inferencia de nuevo.

\subsubsection{14.3 Cómo limpiar el historial}

Haz clic en \textbf{'Limpiar'} (con icono de papelera) en la parte superior del panel de historial. Esto eliminará todos los ejercicios guardados.

\subsubsection{14.4 Indicadores de color en el historial}

Color de etiqueta --- Significado
Verde (\texttt{Valida}) --- La inferencia demostrada era válida
Índigo (\texttt{Valida (Ind.)}) --- Válida pero requiere método indirecto
Rojo (\texttt{Invalida}) --- La inferencia era inválida

\hrulefill

\subsection{15. Los 8 Operadores Lógicos: Referencia Completa}

\subsubsection{15.1 Negación ($ \lnot $) --- NOT}

\begin{itemize}[leftmargin=*]
\item \textbf{Símbolo:} $ \lnot $
\item \textbf{Palabra clave:} NO
\item \textbf{Tipo:} Unario (solo afecta a una variable)
\item \textbf{Significado:} Invierte el valor de verdad. Si P es verdadero, $ \lnot $P es falso.
\item \textbf{Ejemplo:} \texttt{P} = 'No es el caso de P'
\end{itemize}

\subsubsection{15.2 Conjunción ($ \land $) --- AND}

\begin{itemize}[leftmargin=*]
\item \textbf{Símbolo:} $ \land $
\item \textbf{Palabra clave:} Y
\item \textbf{Tipo:} Binario (conecta dos variables)
\item \textbf{Significado:} Solo es verdadero cuando \textbf{ambas} partes son verdaderas.
\item \textbf{Ejemplo:} \texttt{P  Q} = 'P y Q'
\end{itemize}

\subsubsection{15.3 Disyunción ($ \lor $) --- OR}

\begin{itemize}[leftmargin=*]
\item \textbf{Símbolo:} $ \lor $
\item \textbf{Palabra clave:} O
\item \textbf{Tipo:} Binario
\item \textbf{Significado:} Es verdadero cuando \textbf{al menos una} de las partes es verdadera.
\item \textbf{Ejemplo:} \texttt{P  Q} = 'P o Q' (o ambas)
\end{itemize}

\subsubsection{15.4 Disyunción Exclusiva () --- XOR}

\begin{itemize}[leftmargin=*]
\item \textbf{Símbolo:}
\item \textbf{Palabra clave:} O\_EXCLUSIVA
\item \textbf{Tipo:} Binario
\item \textbf{Significado:} Es verdadero cuando \textbf{exactamente una} de las partes es verdadera, pero no ambas.
\item \textbf{Ejemplo:} \texttt{P  Q} = 'P o Q, pero no ambos'
\end{itemize}

\subsubsection{15.5 Implicación ($ \to $) --- IF...THEN}

\begin{itemize}[leftmargin=*]
\item \textbf{Símbolo:} $ \to $
\item \textbf{Palabra clave:} ENTONCES
\item \textbf{Tipo:} Binario
\item \textbf{Significado:} Es falso \textbf{solo cuando} la primera parte es verdadera y la segunda es falsa.
\item \textbf{Ejemplo:} \texttt{P  Q} = 'Si P entonces Q'
\item \textbf{Nota:} Si P es falso, la implicación es verdadera sin importar Q.
\end{itemize}

\subsubsection{15.6 Bicondicional ($ \leftrightarrow $) --- IF AND ONLY IF}

\begin{itemize}[leftmargin=*]
\item \textbf{Símbolo:} $ \leftrightarrow $
\item \textbf{Palabra clave:} SI\_Y\_SOLO\_SI
\item \textbf{Tipo:} Binario
\item \textbf{Significado:} Es verdadero cuando \textbf{ambas partes tienen el mismo valor} (ambas verdaderas o ambas falsas).
\item \textbf{Ejemplo:} \texttt{P  Q} = 'P si y solo si Q'
\end{itemize}

\subsubsection{15.7 Nor ($ \downarrow $) --- NOR}

\begin{itemize}[leftmargin=*]
\item \textbf{Símbolo:} $ \downarrow $
\item \textbf{Palabra clave:} NI
\item \textbf{Tipo:} Binario
\item \textbf{Significado:} Es verdadero \textbf{solo cuando ambas partes son falsas}. Equivale a $ \lnot $(P $ \lor $ Q).
\item \textbf{Ejemplo:} \texttt{P  Q} = 'Ni P ni Q'
\end{itemize}

\subsubsection{15.8 Nand ($ \uparrow $) --- NAND}

\begin{itemize}[leftmargin=*]
\item \textbf{Símbolo:} $ \uparrow $
\item \textbf{Palabra clave:} INCOMPATIBLE
\item \textbf{Tipo:} Binario
\item \textbf{Significado:} Es falso \textbf{solo cuando ambas partes son verdaderas}. Equivale a $ \lnot $(P $ \land $ Q).
\item \textbf{Ejemplo:} \texttt{P  Q} = 'P y Q son incompatibles'
\end{itemize}

\hrulefill

\subsection{16. Las 11 Reglas de Inferencia}

El motor lógico implementa \textbf{11 reglas de inferencia} que usa automáticamente para construir la demostración:

\subsubsection{16.1 Modus Ponendo Ponens (MPP)}
\begin{itemize}[leftmargin=*]
\item \textbf{Si:} P $ \to $ Q y P
\item \textbf{Entonces:} Q
\item \textbf{En palabras:} Si la implicación es verdadera y el antecedente es verdadero, el consecuente es verdadero.
\end{itemize}

\subsubsection{16.2 Modus Tollendo Tollens (MTT)}
\begin{itemize}[leftmargin=*]
\item \textbf{Si:} P $ \to $ Q y $ \lnot $Q
\item \textbf{Entonces:} $ \lnot $P
\item \textbf{En palabras:} Si la implicación es verdadera y el consecuente es falso, el antecedente es falso.
\end{itemize}

\subsubsection{16.3 Silogismo Disyuntivo (SD)}
\begin{itemize}[leftmargin=*]
\item \textbf{Si:} P $ \lor $ Q y $ \lnot $P
\item \textbf{Entonces:} Q
\item \textbf{En palabras:} Si una de dos opciones es verdadera y una es falsa, la otra es verdadera.
\end{itemize}

\subsubsection{16.4 Silogismo Disyuntivo Exclusivo (SDE)}
\begin{itemize}[leftmargin=*]
\item \textbf{Si:} P  Q y P
\item \textbf{Entonces:} $ \lnot $Q
\item \textbf{En palabras:} Si exactamente una de dos opciones es verdadera y una es verdadera, la otra es falsa.
\end{itemize}

\subsubsection{16.5 Silogismo Hipotético (SH)}
\begin{itemize}[leftmargin=*]
\item \textbf{Si:} P $ \to $ Q y Q $ \to $ R
\item \textbf{Entonces:} P $ \to $ R
\item \textbf{En palabras:} Si de P se sigue Q, y de Q se sigue R, entonces de P se sigue R.
\end{itemize}

\subsubsection{16.6 Simplificación}
\begin{itemize}[leftmargin=*]
\item \textbf{Si:} P $ \land $ Q
\item \textbf{Entonces:} P (o Q)
\item \textbf{En palabras:} De una conjunción se puede extraer cualquiera de sus partes.
\end{itemize}

\subsubsection{16.7 Doble Negación}
\begin{itemize}[leftmargin=*]
\item \textbf{Si:} $ \lnot $$ \lnot $P
\item \textbf{Entonces:} P
\item \textbf{En palabras:} Negar dos veces un proposition equivale a la proposición original.
\end{itemize}

\subsubsection{16.8 Conjunción}
\begin{itemize}[leftmargin=*]
\item \textbf{Si:} P y Q (por separado)
\item \textbf{Entonces:} P $ \land $ Q
\item \textbf{En palabras:} Si dos proposiciones son verdaderas, su conjunción también lo es.
\end{itemize}

\subsubsection{16.9 Dilema Constructivo}
\begin{itemize}[leftmargin=*]
\item \textbf{Si:} P $ \to $ Q, R $ \to $ S y P $ \lor $ R
\item \textbf{Entonces:} Q $ \lor $ S
\item \textbf{En palabras:} Si de P se sigue Q, de R se sigue S, y una de las dos (P o R) es verdadera, entonces una de las dos (Q o S) es verdadera.
\end{itemize}

\subsubsection{16.10 Eliminación Bicondicional}
\begin{itemize}[leftmargin=*]
\item \textbf{Si:} P $ \leftrightarrow $ Q y P
\item \textbf{Entonces:} Q
\item \textbf{En palabras:} Si P es equivalente a Q y P es verdadero, entonces Q es verdadero.
\end{itemize}

\subsubsection{16.11 Modus Ponens Bicondicional}
\begin{itemize}[leftmargin=*]
\item \textbf{Si:} P $ \leftrightarrow $ Q y Q
\item \textbf{Entonces:} P
\item \textbf{En palabras:} Si P es equivalente a Q y Q es verdadero, entonces P es verdadero.
\end{itemize}

\hrulefill

\subsection{17. Preguntas Frecuentes (FAQ)}

\subsubsection{¿Qué pasa si escribo una fórmula con errores de sintaxis?}

El sistema muestra un \textbf{aviso de advertencia} en la pestaña del Árbol, indicando qué proposiciones contienen errores. Los errores comunes son:
\begin{itemize}[leftmargin=*]
\item Paréntesis desbalanceados: \texttt{(P  Q} (falta \texttt{)})
\item Operadores consecutivos sin operandos: \texttt{P   Q}
\item Operador al inicio sin operando izquierdo: \texttt{ P Q}
\end{itemize}

Corrige la fórmula y vuelve a intentar.

\subsubsection{¿Puedo usar minúsculas?}

El sistema es \textbf{insensible a mayúsculas} internamente, pero se recomienda usar mayúsculas para las variables (P, Q, R) para mantener consistencia con la notación lógica estándar.

\subsubsection{¿Qué significado tiene 'Método indirecto requerido'?}

Significa que el argumento es válido, pero el motor de demostración directa (forward chaining) no pudo encontrar una derivación paso a paso. Esto ocurre con argumentos que requieren técnicas como:
\begin{itemize}[leftmargin=*]
\item \textbf{Reducción al absurdo:} asumir la negación de la conclusión y llegar a una contradicción.
\item \textbf{Prueba condicional:} asumir una hipótesis temporal para demostrar una implicación.
\end{itemize}

El sistema reconoce que el argumento es válido (no hay contraejemplo) pero no puede mostrar la demostración completa.

\subsubsection{¿Puedo usar más de 4 variables?}

Sí. El teclado virtual ofrece 8 variables (P, Q, R, S, A, B, C, D), pero puedes escribir cualquier letra mayúscula como variable directamente en el campo de premisas.

\subsubsection{¿Qué pasa con los operadores $ \downarrow $ (Nor) y $ \uparrow $ (Nand)?}

Aunque están disponibles en el teclado virtual y el árbol los muestra con sus colores correspondientes, el motor de demostración automática no tiene reglas de inferencia dedicadas para estos operadores. Si tu argumento solo usa $ \downarrow $ o $ \uparrow $, puede que el sistema lo clasifique como 'no demostrable directamente' aunque sea válido.

\subsubsection{¿Cómo funciona el zoom del árbol?}

En la parte inferior del árbol hay controles de zoom:
\begin{itemize}[leftmargin=*]
\item **** (Zoom Out): Reduce el tamaño del árbol (mínimo 60\%).
\item \textbf{+} (Zoom In): Aumenta el tamaño del árbol (máximo 160\%).
\item **** (Restablecer): Vuelve al tamaño normal (100\%).
\end{itemize}

\subsubsection{¿Cómo ver el árbol en pantalla completa?}

Haz clic en el botón \textbf{'Ver a pantalla completa'} debajo del árbol. Se abrirá un modal con el árbol expandido. Para cerrarlo:
\begin{itemize}[leftmargin=*]
\item Haz clic en el botón **** (esquina superior derecha).
\item Haz clic fuera del panel (en el fondo oscuro).
\item Presiona la tecla \textbf{Escape} en tu teclado.
\end{itemize}

\subsubsection{¿Se guarda mi trabajo automáticamente?}

El \textbf{historial} de ejercicios se guarda automáticamente en el navegador (localStorage). Si cierras el navegador y lo vuelves a abrir, los últimos 6 ejercicios seguirán disponibles. Sin embargo, esto es solo en tu navegador; no se sincroniza con otros dispositivos.

\subsubsection{¿Puedo exportar un argumento inválido?}

No. La barra de exportación (Markdown y LaTeX) solo aparece cuando la inferencia es \textbf{válida}. Para un argumento inválido, el panel muestra el diagnóstico y contraejemplo, pero no una demostración exportable.

\hrulefill

\subsection{18. Glosario de Términos}

Término --- Definición simple
\textbf{Proposición} --- Una oración que puede ser verdadera o falsa (ej. 'Llueve').
\textbf{Variable proposicional} --- Una letra que representa una proposición (P, Q, R...).
\textbf{Premisa} --- Una afirmación de la cual se parte para razonar.
\textbf{Conclusión} --- Lo que se quiere demostrar a partir de las premisas.
\textbf{Conectivo lógico} --- Un operador que conecta proposiciones ($ \land $, $ \lor $, $ \to $, $ \lnot $...).
\textbf{Inferencia} --- El proceso de derivar una conclusión a partir de premisas.
\textbf{Argumento} --- Un conjunto de premisas más una conclusión.
\textbf{Argumento válido} --- Un argumento donde la conclusión se sigue necesariamente de las premisas.
\textbf{Argumento inválido} --- Un argumento donde la conclusión NO se sigue de las premisas.
\textbf{Contraejemplo} --- Una asignación de valores que muestra que el argumento es inválido.
\textbf{Falacia} --- Un error de razonamiento que parece válido pero no lo es.
\textbf{AST (Árbol de Sintaxis Abstracta)} --- Representación visual de la estructura jerárquica de una fórmula.
\textbf{Trazabilidad} --- El registro paso a paso de cómo se derivó la conclusión.
\textbf{Modus Ponens} --- Regla: de P $ \to $ Q y P, se deduce Q.
\textbf{Modus Tollens} --- Regla: de P $ \to $ Q y $ \lnot $Q, se deduce $ \lnot $P.
\textbf{Silogismo} --- Regla que encadena implicaciones o disyunciones.
\textbf{Demostración formal} --- Una secuencia de pasos lógicos que prueba la validez de un argumento.
\textbf{Forward chaining} --- Estrategia de demostración que aplica reglas 'hacia adelante' desde las premisas.
\textbf{LaTeX} --- Sistema de tipografía usado en documentos académicos de matemáticas.
\textbf{Markdown} --- Formato simple de escritura para documentos de texto.

\hrulefill

\subsection{Créditos}

Este módulo fue diseñado e implementado por el equipo \textbf{Hijos de Linus} cuyos integrantes estan conformados por:

\begin{itemize}[leftmargin=*]
\item Altamirano Astupiñan Renato
\item Centurión Llatas Alex
\item Espinoza Orrego Piero Arom
\item Mio Caballero Arnold André
\item Morocho Lumbre Juan
\end{itemize}

\hrulefill

> \textbf{Nota para desarrolladores:} Este manual está dirigido a usuarios finales. Si necesitas información técnica sobre la arquitectura del código, consulta los archivos \texttt{ANALISIS\_AST.md} e \texttt{INFORME\_MEJORAS.md} en esta misma carpeta.


\section{Cap\'itulo 3 --- Modus Innova: Cuantificadores}
\section{CAPÍTULO IV: MANUAL DE USUARIO DIDÁCTICO PASO A PASO (''MANUAL PARA TODOS'')}

En cumplimiento riguroso con la segunda obligación exigida por el curso
(''documentar explícitamente de cómo usarlo / manual didáctico para
todos''), este capítulo presenta la guía oficial de usuario para operar
el software QuantifiWeb v3.5. Esta guía permite que tanto alumnos de
primeros ciclos como docentes puedan evaluar cuantificadores y
demostraciones formales sin conocimiento técnico previo.

\subsection{4.1 Visión General de la Interfaz y Temática Neón}

La interfaz de QuantifiWeb v3.5 adopta un estilo visual moderno 'Dark
Glassmorphism' con colores neón adaptativos (Azul \#3B82F6, Cyan \#06B6D4,
Verde \#10B981 y Púrpura \#6366F1).

Detalle de uso adicional para 4.1 Visión General de la Interfaz y
Temática Neón: Esta sección garantiza que cualquier usuario pueda seguir
el flujo de operaciones sin cometer errores sintácticos.

En la parte superior se ubica la barra de navegación que permite
alternar entre las dos herramientas principales: la Pestaña 1 '$ \forall $$ \exists $
Cuantificadores (D $ \subset $ $\mathbb{Z}$)' y la Pestaña 2 '$ \leftrightarrow $ Distribución \& Leyes ($ \Delta $)'.

Detalle de uso adicional para 4.1 Visión General de la Interfaz y
Temática Neón: Esta sección garantiza que cualquier usuario pueda seguir
el flujo de operaciones sin cometer errores sintácticos.

\medskip\hrulefill\medskip
 MANUAL DE USUARIO: 4.1 Visión General de la Interfaz y Temática Neón
Instrucciones operativas para 4.1 Visión General de la Interfaz y
Temática Neón. Guía didáctica para alumnos y profesores.
\medskip\hrulefill\medskip

\medskip\hrulefill\medskip

\subsection{4.2 Guía Paso a Paso para la Evaluación de Cuantificadores Lógicos (D $ \subset $ $\mathbb{Z}$)}

Paso 1: Abra la Pestaña 1 '$ \forall $$ \exists $ Cuantificadores (D $ \subset $ $\mathbb{Z}$)'. En el
desplegable 'Tipo de Cuantificador', seleccione '$ \forall $ Universal' (si desea
probar que la condición se cumple para absolutamente todos los
elementos) o '$ \exists $ Existencial' (si busca la existencia de al menos un
elemento testigo).

Detalle de uso adicional para 4.2 Guía Paso a Paso para la Evaluación de
Cuantificadores Lógicos (D $ \subset $ $\mathbb{Z}$): Esta sección garantiza que cualquier
usuario pueda seguir el flujo de operaciones sin cometer errores
sintácticos.

Paso 2: En la casilla 'Dominio de Discurso D', escriba los números
enteros separados por comas (ejemplo: 1, 2, 3, 4) o especifique rangos
enteros (ejemplo: 0 < x < 90). El sistema validará automáticamente que
no ingrese números decimales. Si ingresa un decimal como 3.5, aparecerá
una alerta de restricción.

Detalle de uso adicional para 4.2 Guía Paso a Paso para la Evaluación de
Cuantificadores Lógicos (D $ \subset $ $\mathbb{Z}$): Esta sección garantiza que cualquier
usuario pueda seguir el flujo de operaciones sin cometer errores
sintácticos.

Paso 3: En la casilla 'Predicado P(x)', escriba la condición matemática
a evaluar (ejemplo: x > 3, x \% 2 = 0). Puede usar comparadores (=, ==,
!=, <>, >, >=, <, <=), operaciones aritmeticas (+, -, *, /, \%, \textasciicircum{}) y
conectores lógicos (\&\&, ||, $ \land $, $ \lor $, y, o).

Detalle de uso adicional para 4.2 Guía Paso a Paso para la Evaluación de
Cuantificadores Lógicos (D $ \subset $ $\mathbb{Z}$): Esta sección garantiza que cualquier
usuario pueda seguir el flujo de operaciones sin cometer errores
sintácticos.

Paso 4: Presione el botón azul ' Evaluar Cuantificador en D'. La
pantalla desplegará la tarjeta de resultado en Verde (si es Verdadero) o
en Rojo (si es Falso), mostrando el elemento testigo x\_w o el
contraejemplo x\_k.

Detalle de uso adicional para 4.2 Guía Paso a Paso para la Evaluación de
Cuantificadores Lógicos (D $ \subset $ $\mathbb{Z}$): Esta sección garantiza que cualquier
usuario pueda seguir el flujo de operaciones sin cometer errores
sintácticos.

Paso 5: Observe la Tabla de Evaluación por Elemento. Cada fila mostrará
el valor de x, la sustitución exacta en el predicado y su resultado
booleano individual.

Detalle de uso adicional para 4.2 Guía Paso a Paso para la Evaluación de
Cuantificadores Lógicos (D $ \subset $ $\mathbb{Z}$): Esta sección garantiza que cualquier
usuario pueda seguir el flujo de operaciones sin cometer errores
sintácticos.

Paso 6: En la parte inferior de la tarjeta, revise el recuadro 'Ley de
De Morgan Equivalente'. El sistema construirá automáticamente la
negación analítica $ \lnot $($ \forall $x P(x)) $ \equiv $ $ \exists $x $ \lnot $P(x), invirtiendo analíticamente la
condición interna.

Detalle de uso adicional para 4.2 Guía Paso a Paso para la Evaluación de
Cuantificadores Lógicos (D $ \subset $ $\mathbb{Z}$): Esta sección garantiza que cualquier
usuario pueda seguir el flujo de operaciones sin cometer errores
sintácticos.

\medskip\hrulefill\medskip
 MANUAL DE USUARIO: 4.2 Guía Paso a Paso para la Evaluación de
Cuantificadores Lógicos (D $ \subset $ $\mathbb{Z}$)
Instrucciones operativas para 4.2 Guía Paso a Paso para la Evaluación
de Cuantificadores Lógicos (D $ \subset $ $\mathbb{Z}$). Guía didáctica para alumnos y
profesores.
\medskip\hrulefill\medskip

\medskip\hrulefill\medskip

\subsection{4.3 Guía Paso a Paso para Demostraciones Paso a Paso con Disyunción Fuerte ($ \Delta $)}

Paso 1: Diríjase a la Pestaña 2 '$ \leftrightarrow $ Distribución \& Leyes ($ \Delta $)'.

Detalle de uso adicional para 4.3 Guía Paso a Paso para Demostraciones
Paso a Paso con Disyunción Fuerte ($ \Delta $): Esta sección garantiza que
cualquier usuario pueda seguir el flujo de operaciones sin cometer
errores sintácticos.

Paso 2: En la casilla 'Fórmula P', escriba su proposición simbólica
formal (ejemplo: ($ \forall $x)\textasciitilde{}[(Px $ \lor $ Mx) $ \to $ Rz] o usando el conector de
Disyunción Fuerte Px $ \Delta $ Qx).

Detalle de uso adicional para 4.3 Guía Paso a Paso para Demostraciones
Paso a Paso con Disyunción Fuerte ($ \Delta $): Esta sección garantiza que
cualquier usuario pueda seguir el flujo de operaciones sin cometer
errores sintácticos.

Paso 3: Utilice el Teclado Virtual de Símbolos en pantalla para insertar
caracteres especiales ($ \forall $, $ \exists $, \textasciitilde{}, $ \land $, $ \lor $, $ \Delta $, $ \to $, $ \leftrightarrow $, $ \equiv $, P, Q, R, S, M, x, y,
z). La tecla $ \Delta $ está destacada en color cyan.

Detalle de uso adicional para 4.3 Guía Paso a Paso para Demostraciones
Paso a Paso con Disyunción Fuerte ($ \Delta $): Esta sección garantiza que
cualquier usuario pueda seguir el flujo de operaciones sin cometer
errores sintácticos.

Paso 4: Presione el botón ' Resolver Paso a Paso (Simplificado)'.

Detalle de uso adicional para 4.3 Guía Paso a Paso para Demostraciones
Paso a Paso con Disyunción Fuerte ($ \Delta $): Esta sección garantiza que
cualquier usuario pueda seguir el flujo de operaciones sin cometer
errores sintácticos.

Paso 5: Examine el desarrollo formal generado. El motor aplicará
secuencialmente las reglas analíticas (Definición de Bicondicional,
Implicación, De Morgan, Disyunción Fuerte y Distribución).

Detalle de uso adicional para 4.3 Guía Paso a Paso para Demostraciones
Paso a Paso con Disyunción Fuerte ($ \Delta $): Esta sección garantiza que
cualquier usuario pueda seguir el flujo de operaciones sin cometer
errores sintácticos.

Paso 6: Observe el Resaltador Neón Cyan. En cada renglón de la
demostración, la subexpresión exacta que sufrió una transformación se
encerrará en una tarjeta con resplandor neón <span
class='highlight-change'>...</span> para que identifique al instante qué
cambió.

Detalle de uso adicional para 4.3 Guía Paso a Paso para Demostraciones
Paso a Paso con Disyunción Fuerte ($ \Delta $): Esta sección garantiza que
cualquier usuario pueda seguir el flujo de operaciones sin cometer
errores sintácticos.

\medskip\hrulefill\medskip
 MANUAL DE USUARIO: 4.3 Guía Paso a Paso para Demostraciones Paso a
Paso con Disyunción Fuerte ($ \Delta $)
Instrucciones operativas para 4.3 Guía Paso a Paso para Demostraciones
Paso a Paso con Disyunción Fuerte ($ \Delta $). Guía didáctica para alumnos y
profesores.
\medskip\hrulefill\medskip

\medskip\hrulefill\medskip

\subsection{4.4 Guía de Resolución de Problemas y Mensajes de Error (Troubleshooting)}

Error 1 'El dominio D opera únicamente sobre números enteros ($\mathbb{Z}$)':
Ocurre cuando se ingresan números con punto decimal (ej. 3.5). Solución:
Reemplace los decimales por valores enteros discretos.

Detalle de uso adicional para 4.4 Guía de Resolución de Problemas y
Mensajes de Error (Troubleshooting): Esta sección garantiza que
cualquier usuario pueda seguir el flujo de operaciones sin cometer
errores sintácticos.

Error 2 'Falta un paréntesis de cierre )': Ocurre cuando los paréntesis
de la expresión no están balanceados. Solución: Verifique que cada
paréntesis abierto '(' tenga su correspondiente cierre ')'.

Detalle de uso adicional para 4.4 Guía de Resolución de Problemas y
Mensajes de Error (Troubleshooting): Esta sección garantiza que
cualquier usuario pueda seguir el flujo de operaciones sin cometer
errores sintácticos.

Error 3 'División entre cero': Ocurre si la sustitución de un elemento x
en el predicado produce un divisor nulo (ej. 5 / (x - 2) cuando x = 2).
Solución: Modifique el dominio para excluir los puntos de
indeterminación.

Detalle de uso adicional para 4.4 Guía de Resolución de Problemas y
Mensajes de Error (Troubleshooting): Esta sección garantiza que
cualquier usuario pueda seguir el flujo de operaciones sin cometer
errores sintácticos.

Error 4 'No hay más pasos aplicables': Ocurre en el motor de leyes
cuando la fórmula alcanza su forma simplificada irreducible.

Detalle de uso adicional para 4.4 Guía de Resolución de Problemas y
Mensajes de Error (Troubleshooting): Esta sección garantiza que
cualquier usuario pueda seguir el flujo de operaciones sin cometer
errores sintácticos.

\medskip\hrulefill\medskip
 MANUAL DE USUARIO: 4.4 Guía de Resolución de Problemas y Mensajes de
Error (Troubleshooting)
Instrucciones operativas para 4.4 Guía de Resolución de Problemas y
Mensajes de Error (Troubleshooting). Guía didáctica para alumnos y
profesores.
\medskip\hrulefill\medskip

\medskip\hrulefill\medskip

\subsection{4.5 Preguntas Frecuentes (FAQ)}

P1: ¿Por qué la Ley de Implicación transforma A $ \to $ B en \textasciitilde{}A $ \lor $ B sin
paréntesis innecesarios? R: El algoritmo refactorizado aísla el
antecedente A y aplica la negación directa, evitando parentizados
redundantes \textasciitilde{}((...)) que confundían a los estudiantes.

Detalle de uso adicional para 4.5 Preguntas Frecuentes (FAQ): Esta
sección garantiza que cualquier usuario pueda seguir el flujo de
operaciones sin cometer errores sintácticos.

P2: ¿Cómo funciona el simplificador de pasos en demostraciones extensas?
R: Si la misma regla (ej. Distribución de Cuantificadores) se aplica
consecutivamente en el mismo nivel, el motor agrupa las transformaciones
en un solo paso elegante.

Detalle de uso adicional para 4.5 Preguntas Frecuentes (FAQ): Esta
sección garantiza que cualquier usuario pueda seguir el flujo de
operaciones sin cometer errores sintácticos.

P3: ¿Qué diferencia a la Disyunción Fuerte ($ \Delta $) de la disyunción normal
($ \lor $)? R: La disyunción inclusiva $ \lor $ permite que ambas proposiciones sean
verdaderas. La disyunción fuerte $ \Delta $ exige que sea exactamente una u otra,
pero NO ambas: A $ \Delta $ B $ \equiv $ (A $ \lor $ B) $ \land $ \textasciitilde{}(A $ \land $ B).

Detalle de uso adicional para 4.5 Preguntas Frecuentes (FAQ): Esta
sección garantiza que cualquier usuario pueda seguir el flujo de
operaciones sin cometer errores sintácticos.

\medskip\hrulefill\medskip
 MANUAL DE USUARIO: 4.5 Preguntas Frecuentes (FAQ)
Instrucciones operativas para 4.5 Preguntas Frecuentes (FAQ). Guía
didáctica para alumnos y profesores.
\medskip\hrulefill\medskip

\medskip\hrulefill\medskip

\section{Cap\'itulo 4 --- Linus: Conjuntos}
\section{Manual de Usuario}

Este manual explica paso a paso cómo instalar, ejecutar y usar la aplicación. Está escrito de forma clara para usuarios sin experiencia en programación. Sigue cada paso en orden. Si algo falla, ve a la sección 8.7 ``Resolución de problemas''.

\subsection{Requisitos previos}
Antes de comenzar, instala el siguiente software en tu computadora:

\paragraph{Node.js (versión 18 LTS o superior):}
Node.js es el entorno de ejecución de JavaScript que necesita la aplicación para funcionar en desarrollo. La versión LTS (Long Term Support) es la más estable.
\begin{enumerate}[noitemsep]
    \item Abre tu navegador y ve a: \url{https://nodejs.org/es/}
    \item Haz clic en el botón verde grande que dice ``LTS''.
    \item Descarga el instalador (\texttt{.msi} en Windows, \texttt{.pkg} en Mac).
    \item Ejecuta el instalador y acepta todas las opciones por defecto.
    \item Cuando termine, cierra y vuelve a abrir tu terminal.
\end{enumerate}

Cómo verificar la instalación: abre una terminal y escribe:
\begin{lstlisting}[language=bash]
node --version # Debe mostrar v18.x.x o superior
npm --version  # Debe mostrar 9.x.x o superior
\end{lstlisting}

\paragraph{Git (sistema de control de versiones):}
Git permite descargar el código del proyecto desde GitHub.
\begin{enumerate}[noitemsep]
    \item Ve a: \url{https://git-scm.com/downloads}
    \item Haz clic en el ícono de tu sistema operativo (Windows/Mac/Linux).
    \item Descarga y ejecuta el instalador.
    \item En Windows: acepta todas las opciones por defecto (incluye ``Git Bash'').
\end{enumerate}

Cómo verificar:
\begin{lstlisting}[language=bash]
git --version # Debe mostrar git version 2.x.x
\end{lstlisting}

\paragraph{Un navegador web moderno:}
La aplicación funciona en cualquier navegador moderno:
\begin{itemize}[noitemsep]
    \item Google Chrome (recomendado) - versión 90 o superior
    \item Mozilla Firefox - versión 88 o superior
    \item Microsoft Edge - versión 90 o superior
    \item Safari - versión 14 o superior (Mac/iOS)
\end{itemize}
No es necesario instalar ningún plugin o extensión adicional.

\subsection{Instalación del proyecto paso a paso}
\paragraph{Paso 1 - Abrir la terminal:}
La terminal (también llamada consola o línea de comandos) es el programa donde escribimos los comandos de instalación.
\begin{itemize}[noitemsep]
    \item En Windows: Presiona la tecla Windows + R, escribe ``cmd'' y presiona Enter; o busca ``PowerShell'' en el menú de inicio y ábrelo.
    \item En Mac: Abre Finder $>$ Aplicaciones $>$ Utilidades $>$ Terminal; o presiona Cmd + Espacio, escribe ``Terminal'' y presiona Enter.
    \item En Linux: Presiona Ctrl + Alt + T.
\end{itemize}

\paragraph{Paso 2 - Descargar el proyecto:}
En la terminal, escribe exactamente este comando y presiona Enter:
\begin{lstlisting}[language=bash]
git clone https://github.com/jordyjosech-prog/logica_simbolica_global.git
\end{lstlisting}
Esto descargará todos los archivos del proyecto a una carpeta llamada \texttt{logica\_simbolica\_global} en tu directorio actual. El proceso puede tardar 30-60 segundos dependiendo de tu conexión a internet. Verás aparecer líneas de texto mientras descarga.

\paragraph{Paso 3 - Entrar a la carpeta del proyecto:}
\begin{lstlisting}[language=bash]
cd logica_simbolica_global
\end{lstlisting}
El comando ``cd'' (change directory) cambia el directorio actual. Después de este comando, todos los comandos siguientes se ejecutarán dentro de la carpeta del proyecto.

\paragraph{Paso 4 - Instalar las dependencias:}
\begin{lstlisting}[language=bash]
npm install
\end{lstlisting}
Este comando lee el archivo \texttt{package.json} y descarga automáticamente todas las librerías que necesita el proyecto (Vue 3, Vite, TypeScript, Vitest, etc.). La primera vez puede tardar entre 1 y 3 minutos. Verás muchas líneas de texto mientras descarga. Al finalizar, aparecerá una línea similar a: \texttt{''added 423 packages in 45s''}.

Si ves advertencias (warnings) en amarillo, son normales y puedes ignorarlas. Si ves errores en rojo, ve a la sección 8.7.

\paragraph{Paso 5 - Ejecutar la aplicación:}
\begin{lstlisting}[language=bash]
npm run dev
\end{lstlisting}
Este comando inicia el servidor de desarrollo de Vite. Verás una salida similar a:
\begin{lstlisting}[language=bash]
  VITE v5.x.x  ready in XXX ms

  -->  Local:   http://localhost:5173/
  -->  Network: use --host to expose
\end{lstlisting}
Copia la URL \texttt{http://localhost:5173/} (o similar) y pégala en tu navegador. La aplicación se abrirá automáticamente. Deja la terminal abierta mientras usas la aplicación; cerrarla detendrá el servidor.

\textit{CONSEJO: Si ya usas el puerto 5173, Vite elegirá automáticamente el siguiente disponible (5174, 5175, etc.). Fíjate en la URL que muestra la terminal.}

\subsection{Uso de la interfaz}
Una vez abierta la aplicación en el navegador, verás cuatro secciones principales:

\subsubsection{Sección 1 - Entrada de datos}
\begin{enumerate}[noitemsep]
    \item Campo ``Universo'': escribe todos los posibles elementos que existen en tu problema, separados por comas. Ejemplo: \texttt{a, b, c, d, e}.
    \item Campo ``Conjunto A'': escribe los elementos del primer conjunto. Deben ser del universo. Ejemplo: \texttt{a, b, c}.
    \item Campo ``Conjunto B'': escribe los elementos del segundo conjunto. Ejemplo: \texttt{b, c, d}.
\end{enumerate}
Puedes usar letras, números o palabras como elementos. Ejemplos válidos: \texttt{''1, 2, 3''} o \texttt{''gato, perro, pez''} o \texttt{''x, y, z''}. Los espacios alrededor de las comas se ignoran automáticamente.

\subsubsection{Sección 2 - Selección de operación}
Desplegable ``Operación'': haz clic para ver las opciones y selecciona la que quieres calcular. Las opciones disponibles son:
\begin{itemize}[noitemsep]
    \item Unión: $A \cup B$ (todos los elementos de A y B juntos)
    \item Intersección: $A \cap B$ (solo los que están en ambos)
    \item Diferencia A-B: $A \setminus B$ (los de A que no están en B)
    \item Diferencia B-A: $B \setminus A$ (los de B que no están en A)
    \item Complemento de A: $U \setminus A$ (los del universo que no están en A)
    \item Complemento de B: $U \setminus B$ (los del universo que no están en B)
    \item Potencia: $\mathcal{P}(A)$ o $\mathcal{P}(\text{resultado})$ - todos los subconjuntos posibles
\end{itemize}
(Solo si elegiste Potencia) Aparece un segundo desplegable para elegir de qué conjunto calcular la potencia.

\subsubsection{Sección 3 - Diagrama de Venn}
El diagrama se actualiza automáticamente al cambiar cualquier dato. No necesitas presionar ningún botón. La región coloreada en el diagrama representa exactamente el resultado de la operación:
\begin{itemize}[noitemsep]
    \item Unión: ambos círculos coloreados
    \item Intersección: solo la zona central (solapada) coloreada
    \item Diferencia A-B: solo la parte izquierda de A (sin la intersección)
    \item Complemento de A: área fuera del círculo A coloreada
    \item Potencia: se muestra el número de subconjuntos en lugar del diagrama
\end{itemize}

\subsubsection{Sección 4 - Resultado}
Debajo del diagrama aparece el resultado en notación matemática de conjuntos:
\begin{itemize}[noitemsep]
    \item \texttt{\{ a, b, c \}} - indica que el resultado contiene los elementos a, b y c.
    \item \texttt{\{\}} - indica que el resultado es el conjunto vacío (sin elementos).
\end{itemize}
Si la operación es Potencia, también aparece el cuadro estadístico con: el valor de $n$ (cantidad de elementos del conjunto base), la fórmula $|\mathcal{P}| = 2^n$ y el número total de subconjuntos.

\subsection{Descripción de operaciones}
\begin{table}[htbp]
\centering
\small
\caption{Ejemplos prácticos de uso de las operaciones}
\begin{tabularx}{\textwidth}{p{2.2cm}p{3.2cm}p{2.8cm}p{3.2cm}X}
\toprule
\textbf{Operación} & \textbf{Entrada} & \textbf{Selector} & \textbf{Resultado esperado} & \textbf{Nota} \\
\midrule
Unión \& A=\{1,2,3\}, B=\{3,4,5\} \& Unión \& \{1, 2, 3, 4, 5\} \& 5 elementos \\
Intersección \& A=\{a,b,c\}, B=\{b,c,d\} \& Intersección \& \{b, c\} \& 2 elementos comunes \\
Diferencia A-B \& A=\{x,y,z\}, B=\{y,z\} \& Diferencia A-B \& \{x\} \& Solo lo de A fuera de B \\
Complemento A \& U=\{1..5\}, A=\{1,2\} \& Complemento de A \& \{3, 4, 5\} \& Fuera de A en U \\
Potencia \& A=\{a,b\} \& Potencia de A \& \{\},\{a\},\{b\},\{a,b\} \& 4 subconjuntos ($2^2$) \\
\bottomrule
\end{tabularx}
\end{table}

\subsection{Exportar resultados}
\begin{itemize}[noitemsep]
    \item \textbf{Botón ``Copiar resultado'':} copia el texto del resultado al portapapeles. Luego puedes pegarlo (Ctrl+V) en Word, Excel, PowerPoint o cualquier otro programa.
    \item \textbf{Captura de pantalla:} En Windows presiona Win+Shift+S para capturar el área que elijas. En Mac presiona Cmd+Shift+4.
    \item \textbf{Impresión:} Presiona Ctrl+P (Windows/Linux) o Cmd+P (Mac) para imprimir la página. Puedes elegir ``Guardar como PDF'' en lugar de imprimir físicamente.
    \item \textbf{Exportar el diagrama SVG:} haz clic derecho sobre el diagrama y elige ``Guardar imagen como...'' para descargar el SVG.
\end{itemize}

\subsection{Despliegue en producción}
\paragraph{Generar el build de producción:}
\begin{lstlisting}[language=bash]
npm run build
\end{lstlisting}
Este comando genera una versión optimizada de la aplicación en la carpeta \texttt{dist/}. El proceso incluye: minificación del JavaScript, eliminación de código no usado (tree-shaking), optimización de assets y generación del \texttt{index.html} final. La carpeta \texttt{dist/} contiene todo lo necesario para publicar la aplicación. No necesitas incluir el resto del proyecto.

\paragraph{Publicar en GitHub Pages (gratis):}
\begin{enumerate}[noitemsep]
    \item Asegúrate de tener el repositorio en GitHub.
    \item En el repositorio, ve a Settings $>$ Pages.
    \item En ``Source'', elige ``GitHub Actions''.
    \item El pipeline \texttt{ci.yml} ya configurado publicará automáticamente al hacer push a main.
    \item La URL pública será: \url{https://[usuario].github.io/[repositorio]/}
\end{enumerate}

\paragraph{Publicar en Netlify (alternativa, más fácil):}
\begin{enumerate}[noitemsep]
    \item Crea una cuenta en \url{https://netlify.com}
    \item Arrastra la carpeta \texttt{dist/} al área de deploy de Netlify.
    \item Netlify genera una URL pública en segundos (p.ej. \url{random-name.netlify.app}).
\end{enumerate}

\subsection{Resolución de problemas comunes}
\paragraph{Problema: La aplicación no carga en el navegador.}
\textbf{Solución:} Verifica que el servidor esté corriendo: en la terminal donde ejecutaste \texttt{''npm run dev''}, no debe haber errores en rojo. Si hay errores, solucionarlos primero. Si la terminal se cerró, vuelve a ejecutar \texttt{npm run dev}.

\paragraph{Problema: npm install falla con errores en rojo.}
\textbf{Solución:} Verifica que Node.js esté instalado con \texttt{''node --version''}. Si el error menciona ``EACCES'' o ``permission denied'', en Mac/Linux ejecuta con sudo: \texttt{''sudo npm install''}. En Windows, abre la terminal como Administrador (clic derecho en PowerShell $>$ Ejecutar como administrador).

\paragraph{Problema: El resultado no se actualiza al escribir.}
\textbf{Solución:} Verifica que los elementos estén separados por comas (no por puntos ni espacios solos). Verifica que el servidor de desarrollo siga corriendo. Prueba refrescar la página con Ctrl+F5.

\paragraph{Problema: El diagrama de Venn no muestra colores correctos.}
\textbf{Solución:} Asegúrate de tener instalada la versión más reciente del proyecto (\texttt{git pull origin main} y \texttt{npm install}). Prueba con un navegador diferente.

\paragraph{Problema: git clone falla con error ``repository not found''.}
\textbf{Solución:} Verifica que tienes acceso al repositorio. Si es privado, necesitas autenticarte con tu cuenta de GitHub: \texttt{git config --global user.email ''tu@email.com''} y \texttt{git config --global user.name ''Tu Nombre''}.

\paragraph{Problema: ``npm'' no se reconoce como comando.}
\textbf{Solución:} Node.js no se instaló correctamente o no está en el PATH del sistema. Desinstala Node.js y reinstálalo. En Windows, asegúrate de marcar ``Add to PATH'' durante la instalación.

\textit{IMPORTANTE: Si ninguna solución funciona, abre un Issue en el repositorio de GitHub describiendo el error exacto (copia y pega el texto del mensaje de error) y tu sistema operativo. El equipo responderá en menos de 48 horas.}

\newpage

\% =========================================================================
\% SECCIÓN 9: PRUEBAS Y VALIDACIÓN
\% =========================================================================
\begin{figure}[H]\centering
\includegraphics[width=0.9\textwidth]{cap_2026-09-10_141523.png}
\caption{Evidencia de interfaz (Equipo Linus)}
\end{figure}
\begin{figure}[H]\centering
\includegraphics[width=0.9\textwidth]{cap_2026-09-10_141540.png}
\caption{Evidencia de interfaz (Equipo Linus)}
\end{figure}
\begin{figure}[H]\centering
\includegraphics[width=0.9\textwidth]{cap_2026-09-10_141604.png}
\caption{Evidencia de interfaz (Equipo Linus)}
\end{figure}
\begin{figure}[H]\centering
\includegraphics[width=0.9\textwidth]{cap_2026-09-10_141618.png}
\caption{Evidencia de interfaz (Equipo Linus)}
\end{figure}
\section{Glosario de términos}
APÉNDICE A --- GLOSARIO DE TÉRMINOS (Tabla 12)
\begin{lstlisting}[language=TypeScript,basicstyle=\ttfamily\scriptsize]
================================================================================
- AST: Abstract Syntax Tree (Arbol de Sintaxis Abstracta).
- Token: Unidad lexica minima del proceso de tokenizacion.
- Parser: Componente que construye el AST verificando la gramatica.
- Tautologia: Formula verdadera bajo cualquier asignacion.
- Contradiccion: Formula falsa bajo cualquier asignacion.
- Contingencia: Formula con valores verdaderos y falsos segun la asignacion.
- Composition API: Estilo de programacion de Vue 3 con ref, computed y watch.
- SFC: Single File Component (.vue).
- Vitest: Framework de pruebas unitarias nativo para Vite.
- Post-orden: Recorrido de arbol que visita hijos antes que el padre.
- Implicacion material: P  Q  ¬P  Q. Falso solo cuando P es V y Q es F.
- Bicondicional: P  Q, verdadero cuando comparten el mismo valor de verdad.

================================================================================
\end{lstlisting}

\section{Procedencia de los cap\'itulos}
Cada cap\'itulo proviene \'integramente de las fuentes, sin reescritura: Cap 1 de \texttt{informe\_sinergia.tex} L339--373 y \texttt{informe\_sinergia.txt} \S8; Cap 2 de \texttt{Instrucciones.md} del equipo Hijos de Linus; Cap 3 del \texttt{.txt} de Modus Innova Cap IV L848--1104; Cap 4 de \texttt{informe\_Linus.tex} L942--1150; Glosario del Ap\'endice A de Sinergia.

\end{document}
